Fix an integer $k\ge 3$.

\begin{itemize}
\item An \emph{additive basis of order $k$} is a set $A\subseteq\mathbb N$ (here $\mathbb N=\{1,2,3,\dots\}$) such that there exists $N_0$ with the property that every integer $n\ge N_0$ can be written as a sum of at most $k$ elements of $A$.
\item A \emph{minimal} additive basis of order $k$ is a basis $B$ of order $k$ such that for every $b\in B$ the set $B\setminus\{b\}$ is \emph{not} a basis of order $k$.
\item For $t\in\{0,1,\dots,k\}$ and $n\in\mathbb Z$, let $r_{A,\le t}(n)$ be the number of representations of $n$ as a sum of at most $t$ elements of $A$, where:
  \begin{itemize}
  \item repetition of summands is allowed;
  \item order of summands is ignored (so we count nondecreasing $t'$-tuples for some $0\le t'\le t$);
  \item we include the empty sum as the unique representation of $0$ (so $r_{A,\le t}(0)\ge 1$).
  \end{itemize}
  Define $r_A(n):=r_{A,\le k}(n)$.
\end{itemize}

(Counting ordered representations instead changes $r_A(n)$ by at most a factor $k!$ for fixed $k$, so the qualitative question ``$\ge c\log n$'' is insensitive to this convention.)
